\chapter{Субстрат: типизированный темпоральный граф}
\label{ch:substrate}

\section{Что такое узел и что такое ребро}

Память Рин — направленный граф $G=(V,E)$, в котором узлы несут содержание
(воспоминания, факты, сущности, концепты), а рёбра — типизированные и знаковые
отношения между ними. Ключевая мысль главы: и узлы, и рёбра несут
\emph{метаданные}, и именно метаданные, а не сам факт связи, делают возможными
все дальнейшие механизмы.

\begin{definition}[Узел памяти]
Узел $v\in V$ — это запись с содержанием $\mathbf{x}_v$ и набором полей состояния,
над которыми работают операторы памяти. Узел одновременно является и «данными»
(что помнится), и «убеждением» (насколько это считается истинным), и «следом»
(насколько активно и свежо).
\end{definition}

\subsection{Схема узла}

\begin{longtable}{p{0.24\textwidth} p{0.10\textwidth} p{0.56\textwidth}}
\toprule
\textbf{Поле} & \textbf{Тип} & \textbf{Смысл и кто с ним работает} \\
\midrule
\endhead
\texttt{id} & int & идентификатор \\
\texttt{content} & text & содержание в естественном виде \\
\texttt{embedding} $\mathbf{x}_v$ & vec$[d]$ & векторное представление; retrieval-релевантность, мост баз \\
\texttt{confidence} $c_v$ & $[0,1]$ & уверенность как в убеждении; пропагация, ревизия \\
\texttt{valence} & vec$[2]$ & валентность–возбуждение; сцепка с настроением \\
\texttt{layer}/\texttt{ts} & enum/time & масштаб (сессия\ldots жизнь) и/или метка времени; консолидация \\
\texttt{provenance} & list & источники и/или обоснования; TMS, ревизия \\
\texttt{need\_to\_*} & flags & «исследовать/пересмотреть»; любопытство, скраб \\
\texttt{activation} $a_v$ & float & транзиентная активация; обход \\
\texttt{degree}/\texttt{centrality} & float & производные структурные; топология \\
\bottomrule
\end{longtable}

\subsection{Схема ребра}

Ребро $(u\to v)$ несёт: направленность (\texttt{src},\texttt{dst}); тип отношения
\texttt{type} (\textsc{isA}, \textsc{partOf}, \textsc{supports},
\textsc{contradicts}, \textsc{associates}, \dots); \textbf{знак} $\pm$ (минимум,
без которого не работает знаковая пропагация); вес \texttt{weight} $w_{uv}$ —
эвиденциальную силу связи.

\begin{intuition}[Веса — про свидетельство, а не про характер]
Забегая вперёд (глава~\ref{ch:personality}): вес ребра выражает
\emph{эвиденциальную} силу связи и уверенность, а \emph{не} «характер». Соблазн
задать личность статическими множителями на рёбра-реакции («грусть$\to$сарказм
$\times 2{,}0$») мы разберём и отвергнем: характер должен входить как
\emph{смещение во время обхода} от текущего состояния, а не как вшитый в граф
множитель.
\end{intuition}

\section{Почему рёбра направленные, типизированные и знаковые}

Плоская модель «узлы и связи одинакового сорта» ломается сразу на трёх вещах:
направленность (асимметрия причинности и времени), тип (нельзя валить covalent-связь,
близость, «подтверждает» и «является частью» в один линейный канал) и знак (без
него не отличить «всё усилилось» от «всё перевернулось»).

Правильный оператор поэтому не общий лапласиан сглаживания, а
\emph{relation-specific} передача сообщения. Сообщение по ребру:
\begin{equation}
m_{u\to v} \;=\; \phi_{r_{uv}}\!\big(h_u,\, h_v,\, e_{uv},\, \Delta t,\, c_{uv},\, s_{uv}\big),
\label{eq:message}
\end{equation}
где $r_{uv}$ — тип отношения, $\Delta t$ — возраст/временная дистанция, $c_{uv}$ —
уверенность, $s_{uv}$ — источник/эпистемический статус. Затем узел агрегирует
входящие сообщения — и здесь важно не одну статистику, а набор (сумма/среднее/макс/
дисперсия по типам отношений), потому что для памяти существенны и кратность
подтверждений, и пропорции, и наличие экстремального свидетельства.

\begin{workedbox}[Почему одна статистика теряет смысл]
Пусть узел «юзер X» имеет соседей-свидетельств: три источника подтверждают факт,
один противоречит. \emph{Среднее} склеит $\{+,+,+,-\}$ и $\{+,-\}$, если пропорции
совпадут, потеряв, что в первом случае \emph{четыре} независимых свидетельства, а
во втором \emph{два}. Для памяти число независимых подтверждений принципиально:
«три источника против одного» — не то же, что «один против одного». Поэтому
агрегатор хранит и кратность (сумма), и пропорции (среднее), и наличие сильного
контрсвидетельства (макс по знаку «против»).
\end{workedbox}

\begin{intuition}[Аналогия с электрической цепью]
Полезно думать о графе как о цепи с проводниками разного сорта. Направленность —
это диоды; тип ребра — материал проводника со своей проводимостью; знак — это
полярность. Плоская модель — цепь, где все элементы заменены одинаковыми
резисторами: она проведёт ток, но потеряет всю логику схемы.
\end{intuition}

\begin{center}
\begin{tikzpicture}[>=Stealth, node distance=22mm,
  every node/.style={font=\small},
  mem/.style={draw, rounded corners, minimum width=15mm, minimum height=7mm}]
  \node[mem] (a) {разбить кружку};
  \node[mem, right=30mm of a] (b) {грусть};
  \node[mem, below=13mm of b] (c) {неуклюжесть};
  \node[mem, below=13mm of a] (d) {юзер: хотел новую};
  \draw[->, thick] (a) -- node[above,font=\scriptsize]{$+$ xReact} (b);
  \draw[->, thick] (a) -- node[left,font=\scriptsize]{$+$ assoc} (c);
  \draw[->, thick, red!60!black] (d) -- node[below,font=\scriptsize]{$-$ contradicts} (b);
\end{tikzpicture}\\[2pt]
{\small Рис.: типизированные знаковые рёбра. Личная память «юзер хотел новую»
\emph{противоречит} базовой реакции «грусть» — это и делает возможным уместное
«поздравляю».}
\end{center}

\section{Провенанс и обоснование}
\label{sec:prov}

Поле \texttt{provenance} шире, чем «источник». Для \emph{исходных} фактов провенанс
— откуда факт узнан. Для узлов, \emph{порождённых выводом} (fusion), провенанс
должен хранить \emph{обоснование} — множество родительских узлов, из которых вывод
получен. Эта деталь окажется критической в главе~\ref{ch:epistemics}: именно
наличие обоснований позволяет, опровергнув посылку, автоматически пометить на
пересмотр её следствия. Без записанных обоснований опровержение посылки не находит
потомков — и рождается «ползучая вера».

\section{Что значит «коммит» и кто его гейтит}

Важнейшая граница схемы — между \emph{активным рабочим убеждением} и \emph{записью}.
Не всё, что всплыло или сгенерировалось, попадает в граф как факт.

\begin{principle}[Гипотеза, а не факт]
Всё, что порождено выводом или генерацией (fusion, реакция, ассоциация), пишется
\emph{провизорно}: с низкой уверенностью, с провенансом-обоснованием и, при
необходимости, с флагом \texttt{need\_to\_discover}. Повышение до «записи» (переход
в более медленный слой) происходит только через накопление подтверждающих
свидетельств. Управление — на границе коммита.
\end{principle}

\begin{lstlisting}[style=pseudo]
def commit(G, item, kind):
    if kind == "observed":          # прямое наблюдение
        c0 = source_trust(item.source)
    if kind == "inferred":          # порождено fusion/выводом
        c0 = LOW                    # гипотеза, не факт
        item.provenance = item.parents      # обоснование для TMS
    node = G.add_node(item, confidence=c0, layer="session")
    propagate(G, node)              # см. главу про лавину (со знаком, карантином)
    return node
\end{lstlisting}

\section{Кража схемы у Wikidata}

Есть готовая, обкатанная модель данных, уже решившая ровно те задачи, что встают
перед метаданными Рин, — модель утверждений \emph{Wikidata}. Брать её
\emph{содержание} (энциклопедические факты) для существа, мимикрирующего под
человека, не нужно: человек помнит мало фактов, его сила в ассоциативности. Но её
\emph{схему} стоит перенять целиком:

\begin{itemize}
\item \textbf{квалификаторы} (точка во времени, начало/конец) — темпоральная
валидность: факт истинен \emph{в интервале}, а не вечно;
\item \textbf{references} — провенанс/обоснование;
\item \textbf{ранг} (\emph{preferred}/\emph{normal}/\emph{deprecated}) —
\emph{недеструктивная ревизия}: устаревшее не удаляют, а депрецируют, помечая новое
предпочтительным; история сохраняется для машинерии поддержки истинности.
\end{itemize}
Иными словами, факт стоит моделировать как
\[
\text{(утверждение)} + \text{(интервал валидности)} + \text{(провенанс)} +
\text{(статус/ранг)},
\]
и это буквально схема Wikidata. Не изобретать метаданные заново, а взять
проверенный шаблон.

\section{Врождённое и приобретённое: не граница, а континуум}

Соблазнительно разделить узлы на «врождённые константы» (эмбриональный граф) и
«приобретённые» (нажитые, с временными метками). Как первое приближение полезно, но
как жёсткая граница неверно. Консолидация \emph{двигает} записи через границу
(эпизод $\to$ обобщённое знание), а ревизия должна уметь \emph{понижать} и
«врождённое» — иначе унаследованную ошибку здравого смысла Рин не исправит никогда.
Правильная картина — \emph{континуум глубины консолидации}, где порог опрокидывания
убеждения растёт с глубиной, но нигде не становится бесконечным. К оси
«пластичность\,--\,упрямство» мы вернёмся не раз (главы~\ref{ch:propagation},
\ref{ch:epistemics}).
